\chapter{Kansrekening op aftelbare ruimten}\label{ch:b2:proba}

De laatste drie hoofdstukken ontwikkelen de moderne kansrekening:
kansmaten op \omterm{def:b2:structures:countable}{aftelbare} \omterm{def:b2:proba:space}{uitkomstenruimten}, discrete
toevalsveranderlijken en \omterm{ex:b2:powerseries:fibonacci}{genererende functies}. De eindige theorie
van het bovenbouwvolume krijgt haar volledige infrastructuur:
$\sigma$-additiviteit vervangt de eindige additiviteit, en de
machinerie van de \omterm{def:b2:series:summable}{sommeerbare families} uit \cref{ch:b2:series} is
precies wat oneindige \omterm{def:b2:proba:space}{uitkomstenruimten} hanteerbaar maakt. De
kernresultaten hier zijn de \omterm{def:b2:metric:continuity}{continuïteit} van de kans langs monotone
rijen \omterm{def:b2:proba:space}{gebeurtenissen} en het lemma van Borel--Cantelli.

\section{Kansruimten}

\begin{definition}[Aftelbare kansruimte]\label{def:b2:proba:space}
Zij $\Omega$ een niet-lege eindige of \omterm{def:b2:structures:countable}{aftelbare verzameling} (de
\emph{uitkomstenruimte}\index{uitkomstenruimte}). Een
\emph{kansmaat}\index{kansmaat} op $\Omega$ is een afbeelding $\P$
van de verzameling $\mathcal{P}(\Omega)$ van alle
deelverzamelingen van $\Omega$
(\emph{gebeurtenissen}\index{gebeurtenis}) naar $[0, 1]$ zodanig
dat:
\begin{enumerate}
\item $\P(\Omega) = 1$;
\item ($\sigma$-additiviteit) voor elke rij $(A_n)_{n\in\N}$
paarsgewijs disjuncte gebeurtenissen geldt
\[
\P\Bigl(\,\bigcup_{n \in \N} A_n\Bigr)
= \sum_{n=0}^{\infty} \P(A_n) .
\]
\end{enumerate}
Het paar $(\Omega, \P)$ is een (\omterm{def:b2:structures:countable}{aftelbare})
\emph{kansruimte}\index{kansruimte}.
\end{definition}

\begin{remark}
Op een \omterm{def:b2:structures:countable}{aftelbare} $\Omega$ mogen wij alle deelverzamelingen als
\omterm{def:b2:proba:space}{gebeurtenissen} nemen; op overaftelbare ruimten (zoals nodig voor
\omterm{def:b2:metric:continuity}{continue} modellen in bachelorjaar 3) kan dat niet meer, en beperkt
men $\P$ tot een geschikte collectie \omterm{def:b2:proba:space}{gebeurtenissen}, een
\emph{$\sigma$-algebra}. Alle formules van dit hoofdstuk overleven
die veralgemening woordelijk.
\end{remark}

\begin{proposition}[Elementaire regels]\label{prop:b2:proba:rules}
Voor \omterm{def:b2:proba:space}{gebeurtenissen} $A, B$ en een \omterm{def:b2:proba:space}{kansmaat} $\P$ geldt:
$\P(\emptyset) = 0$; $\P$ is eindig additief; $\P(A^c) = 1 -
\P(A)$; is $A \subseteq B$, dan $\P(A) \leq \P(B)$; en
\[
\P(A \cup B) = \P(A) + \P(B) - \P(A \cap B) .
\]
\end{proposition}

\begin{proof}
De $\sigma$-additiviteit toepassen op $A_0 = \Omega$, $A_n =
\emptyset$ ($n \geq 1$) geeft $1 = 1 +
\sum_{n\geq1}\P(\emptyset)$, dus $\P(\emptyset) = 0$; een eindige
disjuncte vereniging met lege verzamelingen opvullen geeft dan de
eindige additiviteit. De rest volgt als in het eindige geval
(bovenbouwvolume): $1 = \P(A) + \P(A^c)$ uit $\Omega = A \sqcup
A^c$; $\P(B) = \P(A) + \P(B \setminus A) \geq \P(A)$ wanneer $A
\subseteq B$; en ontbinden in drie disjuncte stukken geeft
\begin{align*}
\P(A \cup B) &= \P(A \setminus B) + \P(B \setminus A) +
\P(A \cap B)\\
&= \bigl(\P(A) - \P(A\cap B)\bigr) + \bigl(\P(B) - \P(A\cap
B)\bigr) + \P(A \cap B),
\end{align*}
wat de in- en uitsluiting is; de algemene versie met $n$
verzamelingen is \cref{exo:b2:proba:4}.
\end{proof}

\begin{proposition}[Verdelingen op een aftelbare
ruimte]\label{prop:b2:proba:distribution}
Een \omterm{def:b2:proba:space}{kansmaat} geven op een \omterm{def:b2:structures:countable}{aftelbare} $\Omega = \{\omega_0,
\omega_1, \dots\}$ komt precies neer op het geven van gewichten
$p_i = \P(\{\omega_i\}) \geq 0$ met $\sum_i p_i = 1$; dan is voor
elke $A \subseteq \Omega$
\[
\P(A) = \sum_{\omega \in A} \P(\{\omega\}) ,
\]
een (\omterm{def:b2:series:def}{absoluut} convergente) deelsom van de familie $(p_i)$.
\end{proposition}

\begin{proof}
Gegeven $\P$ vormen de singletons $\{\omega\}$, $\omega \in A$, een
\omterm{def:b2:structures:countable}{aftelbare} disjuncte overdekking van $A$, dus dwingt de
$\sigma$-additiviteit af dat
\[
\P(A) = \sum_{\omega\in A}\P(\{\omega\}),
\]
een onvoorwaardelijke deelsom van de niet-negatieve \omterm{def:b2:series:summable}{sommeerbare
familie} $(p_i)$ --- herschikken is onschadelijk precies omdat de
termen niet-negatief zijn (\cref{ch:b2:series}); in het bijzonder
is $\sum_ip_i = \P(\Omega) = 1$. Omgekeerd, gegeven niet-negatieve
gewichten met totale som $1$, definieer $\P(A) = \sum_{\omega \in
A}p_\omega$: de familie is \omterm{def:b2:series:summable}{sommeerbaar}, en de
$\sigma$-additiviteit is precies de stelling over sommeren in
pakketten uit \cref{ch:b2:series}, toegepast op de partitie van
$\bigcup A_n$ in de $A_n$.
\end{proof}

\begin{example}[Meetkundig model: wachten op de eerste
kop]\label{ex:b2:proba:geometric}
Werp herhaaldelijk met een munt met kanskop $p \in \intoo{0}{1}$,
en laat $\Omega = \N^* \cup \{\infty\}$ de rang van de eerste kop
registreren. De natuurlijke gewichten zijn
\[
\P(\{k\}) = (1 - p)^{k-1}p
\quad (k \in \N^*),
\qquad
\P(\{\infty\}) = 0 ,
\]
een \omterm{def:b2:proba:space}{kansmaat} omdat $\sum_{k\geq1}(1-p)^{k-1}p = \frac{p}{1 -
(1-p)} = 1$: met kans $1$ eindigt het spel --- maar de
\omterm{def:b2:proba:space}{uitkomstenruimte} moet toch de mogelijkheid bevatten dat het niet
eindigt. De \omterm{def:b2:structures:countable}{aftelbare} additiviteit is wat ons toelaat te beweren
dat $\P(\text{het spel eindigt}) = \sum_k \P(\{k\})$.
\end{example}

\begin{theorem}[Monotone continuïteit]\label{thm:b2:proba:continuity}
Zij $(A_n)$ een rij \omterm{def:b2:proba:space}{gebeurtenissen}.
\begin{enumerate}
\item Is $A_n \subseteq A_{n+1}$ voor alle $n$
(\emph{stijgend}), dan is $\P\bigl(\bigcup_n A_n\bigr) =
\lim_{n\to\infty} \P(A_n)$.
\item Is $A_n \supseteq A_{n+1}$ voor alle $n$
(\emph{dalend}), dan is $\P\bigl(\bigcap_n A_n\bigr) =
\lim_{n\to\infty} \P(A_n)$.
\end{enumerate}
\end{theorem}

\begin{proof}
\emph{1.} Maak disjunct: zij $B_0 = A_0$ en $B_n = A_n \setminus
A_{n-1}$. De $B_n$ zijn paarsgewijs disjunct met $\bigcup_{k \leq
n} B_k = A_n$ en $\bigcup_n B_n = \bigcup_n A_n$. Volgens de
$\sigma$-additiviteit en de eindige additiviteit is
\[
\P\Bigl(\bigcup_n A_n\Bigr)
= \sum_{n=0}^\infty \P(B_n)
= \lim_{N\to\infty}\sum_{n=0}^N \P(B_n)
= \lim_{N\to\infty}\P(A_N) .
\]
\emph{2.} Ga over op de complementen: $(A_n^c)$ is stijgend met
vereniging $\bigl(\bigcap A_n\bigr)^c$; pas deel 1 toe: $1 -
\P(\bigcap A_n) = \lim (1 - \P(A_n))$.
\end{proof}

\begin{corollary}[Aftelbare subadditiviteit]\label{cor:b2:proba:subadd}
Voor elke rij \omterm{def:b2:proba:space}{gebeurtenissen} is $\P\bigl(\bigcup_n A_n\bigr) \leq
\sum_{n=0}^\infty \P(A_n)$.
\end{corollary}

\begin{proof}
De eindige subadditiviteit $\P(A_0 \cup \dots \cup A_N) \leq
\sum_0^N \P(A_n)$ volgt met inductie uit de in- en uitsluiting (of
uit de additiviteit over de disjunct gemaakte $B_n \subseteq
A_n$). Laat $N \to \infty$: het linkerlid convergeert naar
$\P(\bigcup_n A_n)$ wegens de monotone \omterm{def:b2:metric:continuity}{continuïteit} toegepast op de
stijgende rij $C_N = A_0 \cup \dots \cup A_N$.
\end{proof}

\begin{example}[De somgrens: ruw maar onverwoestbaar]\label{ex:b2:proba:unionbound}
De subadditiviteit met eindig veel \omterm{def:b2:proba:space}{gebeurtenissen} --- de
\emph{somgrens} --- ruilt nauwkeurigheid in voor universaliteit.
Voor het verjaardagsprobleem met $23$ mensen geeft de kans op een
botsing begrenzen door de som over de paren
\[
\P(\text{botsing}) \leq \binom{23}2\cdot\frac1{365}
= \frac{253}{365} \approx 0.693 ,
\]
tegenover de werkelijke $0.507$: er ruim naast, omdat botsingen
elkaar overlappen. Toch heeft de grens \emph{geen}
\omterm{def:b2:proba:independence}{onafhankelijkheid} nodig, geen gezamenlijke verdeling, niets dan de
kansen op de paren --- en daarom is in de weekendopgave en overal
in \cref{ch:b2:randomvar} de somgrens het eerst getrokken wapen:
blijkt zij klein te zijn, dan is de zaak beslecht zonder verdere
modellering.
\end{example}

\begin{example}[Een zes komt, uiteindelijk]\label{ex:b2:proba:sixeventually}
Gooi eeuwig met een eerlijke dobbelsteen en zij $B_n = {}$``minstens
één zes onder de eerste $n$ worpen'', een stijgende rij
\omterm{def:b2:proba:space}{gebeurtenissen} met $\P(B_n) = 1 - (5/6)^n$. De monotone
\omterm{def:b2:metric:continuity}{continuïteit} geeft
\[
\P(\text{er verschijnt uiteindelijk een zes})
= \P\Bigl(\bigcup_nB_n\Bigr)
= \lim_n\bigl(1 - (5/6)^n\bigr) = 1 .
\]
Het punt is niet de (voor de hand liggende) limiet maar de logische
stap: ``uiteindelijk'' is een \omterm{def:b2:proba:space}{gebeurtenis} over \emph{oneindig veel}
worpen, buiten het bereik van de eindige additiviteit, en de
monotone \omterm{def:b2:metric:continuity}{continuïteit} --- dat wil zeggen de $\sigma$-additiviteit
--- is precies het axioma dat er een kans aan toekent. Elke uitspraak
``bijna zeker'' in de rest van dit boek gaat door deze ene smalle
deur.
\end{example}

\section{Voorwaardelijkheid en onafhankelijkheid}

\begin{definition}[Voorwaardelijke kans]\label{def:b2:proba:conditional}
Voor \omterm{def:b2:proba:space}{gebeurtenissen} $A, B$ met $\P(B) > 0$ is de
\emph{voorwaardelijke kans}\index{voorwaardelijke kans} van $A$
gegeven $B$
\[
\P(A \mid B) = \frac{\P(A \cap B)}{\P(B)} .
\]
De afbeelding $A \mapsto \P(A \mid B)$ is zelf een \omterm{def:b2:proba:space}{kansmaat} op
$\Omega$.
\end{definition}

\begin{remark}
Dat $A \mapsto \pcond BA$ opnieuw een \omterm{def:b2:proba:space}{kansmaat} is, verdient een
ogenblik: $\pcond B\Omega = 1$ en de $\sigma$-additiviteit gaan
door het quotiënt heen omdat doorsnijden met $B$ disjuncte
verenigingen eerbiedigt. Het praktische gevolg: elke identiteit van
dit hoofdstuk --- in- en uitsluiting, monotone \omterm{def:b2:metric:continuity}{continuïteit},
Borel--Cantelli --- mag \emph{na} het voorwaardelijk stellen worden
toegepast, zonder nieuwe bewijzen. Kansrekenaars ``werken onder
$\pcond B{\cdot}$'' precies om die reden voortdurend.
\end{remark}

\begin{example}[Voorwaardelijk stellen kan uniformiteit
scheppen]\label{ex:b2:proba:sumseven}
Gooi met twee eerlijke dobbelstenen en stel voorwaardelijk dat de
som $7$ is: voor elke $k \in \intint16$ is
\[
\pcond{\{S = 7\}}{X = k}
= \frac{\P(X = k,\ Y = 7 - k)}{\P(S = 7)}
= \frac{1/36}{6/36} = \frac16 :
\]
gegeven een som van $7$ is de eerste dobbelsteen precies uniform
--- $7$ is de enige som die met elk aantal ogen verenigbaar is, dus
wist het voorwaardelijk stellen alle informatie over $X$ uit. Elke
andere som scheeftrekt de verdeling (gegeven $S = 4$ is de eerste
dobbelsteen alleen uniform op $\{1, 2, 3\}$). Een voorwaardelijke
verdeling berekenen betekent de gezamenlijke gewichten langs de
voorwaardelijke \omterm{def:b2:proba:space}{gebeurtenis} hernormaliseren, meer niet.
\end{example}

\begin{example}[De tweede trekking is even goed als de
eerste]\label{ex:b2:proba:seconddraw}
Een urne bevat $3$ witte en $2$ zwarte ballen; trek er twee zonder
teruglegging. Iedereen is het erover eens dat $\P(W_1) = \frac35$;
wat is $\P(W_2)$? Totale kans langs de eerste trekking:
\[
\P(W_2) = \pcond{W_1}{W_2}\,\P(W_1) +
\pcond{B_1}{W_2}\,\P(B_1)
= \frac24\cdot\frac35 + \frac34\cdot\frac25
= \frac{12}{20} = \frac35 :
\]
precies $\P(W_1)$. Er was geen berekening nodig: wegens de
symmetrie heeft elke bal evenveel kans om als tweede te worden
getrokken, dus heeft de tweede trekking --- \emph{onvoorwaardelijk}
--- dezelfde verdeling als de eerste. Voorwaardelijk stellen op de
eerste uitkomst verandert de kansen; haar niet kennen niet. Dit
argument met verwisselbaarheid keert in het volgende hoofdstuk
terug voor het steekproeven zonder teruglegging, waar het het
hypergeometrische gemiddelde $np$ geeft, zonder enige binomiale
identiteit.
\end{example}

\begin{theorem}[Samengestelde kansen, totale kans,
Bayes]\label{thm:b2:proba:bayes}
\leavevmode
\begin{enumerate}
\item (Kettingregel) Is $\P(A_1 \cap \dots \cap A_{n-1}) > 0$, dan
\[
\P(A_1 \cap \dots \cap A_n)
= \P(A_1)\,\P(A_2 \mid A_1)\cdots
\P(A_n \mid A_1 \cap \dots \cap A_{n-1}) .
\]
\item (Totale kans) Is $(B_i)_{i \in I}$ een eindige of \omterm{def:b2:structures:countable}{aftelbare}
partitie van $\Omega$ met $\P(B_i) > 0$, dan is voor elke
\omterm{def:b2:proba:space}{gebeurtenis} $A$
\[
\P(A) = \sum_{i \in I} \P(A \mid B_i)\,\P(B_i) .
\]
\item (Bayes) Onder dezelfde hypothesen geldt, als bovendien
$\P(A) > 0$:
\[
\P(B_j \mid A)
= \frac{\P(A \mid B_j)\,\P(B_j)}
       {\sum_{i \in I} \P(A \mid B_i)\,\P(B_i)} .
\]
\end{enumerate}
\end{theorem}

\begin{proof}
\emph{1.} Schrijf elke \omterm{def:b2:proba:conditional}{voorwaardelijke kans} als een quotiënt: het
rechterlid is
\[
\P(A_1)\cdot\frac{\P(A_1 \cap A_2)}{\P(A_1)}\cdot
\frac{\P(A_1 \cap A_2 \cap A_3)}{\P(A_1 \cap A_2)}\cdots
\frac{\P(A_1 \cap \dots \cap A_n)}{\P(A_1 \cap \dots \cap
A_{n-1})},
\]
een telescoperend product: elke noemer heft de voorgaande teller
op, en $\P(A_1 \cap \dots \cap A_n)$ blijft over. Alle noemers zijn
$\geq \P(A_1 \cap \dots \cap A_{n-1}) > 0$ wegens de monotonie, dus
verdwijnt er niets. (De hypothese bewaakt precies dit:
voorwaardelijk stellen op een \omterm{def:b2:proba:space}{gebeurtenis} met kans nul is niet
gedefinieerd.)
\emph{2.} De verzamelingen $A \cap B_i$ zijn paarsgewijs disjunct
met vereniging $A$; pas de ($\sigma$-)additiviteit en de definitie
van het voorwaardelijk stellen toe.
\emph{3.} Beide leden van $\P(B_j \mid A)\P(A) = \P(A \mid
B_j)\P(B_j)$ zijn gelijk aan $\P(A \cap B_j)$; deel door $\P(A)$ en
werk $\P(A)$ uit met de totale kans.
\end{proof}

\begin{example}[De verjaardagsbotsing, met de
kettingregel]\label{ex:b2:proba:birthday}
Met $n$ mensen wier verjaardagen \omterm{def:b2:proba:independence}{onafhankelijk} en uniform over
$365$ dagen verdeeld zijn, zij $D_n = {}$``alle $n$ verjaardagen
verschillen''. Persoon voor persoon voorwaardelijk stellen
(kettingregel) geeft
\[
\P(D_n) = \prod_{k=1}^{n-1}\Bigl(1 - \frac{k}{365}\Bigr),
\]
waarbij elke nieuwe persoon de $k$ reeds bezette dagen moet
vermijden. Voor $n = 23$: $\P(D_{23}) \approx 0.493$ --- een
gedeelde verjaardag is al waarschijnlijker dan niet. De heuristiek
die de kleinheid van $23$ verklaart: logaritmen nemen geeft $-\ln
\P(D_n) \approx \sum_{k<n}\frac k{365} = \frac{\binom n2}{365}$, en
$\binom{23}2 = 253$ geeft $253/365 \approx 0.693 \approx \ln 2$.
Wat telt is het aantal \emph{paren}, dat kwadratisch groeit:
botsingsproblemen leven op de schaal $n \sim \sqrt{365}$, niet $n
\sim 365$ --- de verjaardagsparadox is een vierkantswortel in
vermomming.
\end{example}

\begin{example}[Monty Hall, met Bayes]\label{ex:b2:proba:montyhall}
Achter een van drie deuren zit uniform verdeeld een prijs. Je kiest
deur $1$; de presentator, die weet waar de prijs is, opent een van
de andere deuren, altijd een lege (uniform kiezend wanneer hij de
keuze heeft), zeg deur $3$. Zij $B_i = {}$``prijs achter deur $i$''
en $A = {}$``de presentator opent deur $3$''. Dan is
$\pcond{B_1}{A} = \frac12$, $\pcond{B_2}{A} = 1$ en
$\pcond{B_3}{A} = 0$, dus geeft Bayes
(\cref{thm:b2:proba:bayes})
\[
\P(B_2 \mid A)
= \frac{1\cdot\frac13}
{\frac12\cdot\frac13 + 1\cdot\frac13 + 0\cdot\frac13}
= \frac23 :
\]
van deur wisselen wint twee keer op drie. De berekening lokaliseert
de populaire verwarring precies: de zet van de presentator is
\emph{informatief} (hij kon deur $2$ niet \omterm{def:b2:metric:topology}{openen} als de prijs daar
zat), en de formule van Bayes is het boekhoudkundige hulpmiddel dat
die asymmetrie in de $\frac23$ omzet. Voorwaardelijk stellen op
``wat werd gezien'' in plaats van op ``wat waar is'' is de hele
kunst van de formule.
\end{example}

\begin{example}[De twee weddenschappen van de Chevalier de
Méré]\label{ex:b2:proba:demere}
Twee weddenschappen uit de zeventiende eeuw, beslecht door
\omterm{def:b2:proba:independence}{onafhankelijkheid}. Weddenschap één: minstens één zes in $4$ worpen
met een dobbelsteen,
\[
\P = 1 - \Bigl(\frac56\Bigr)^{\!4} \approx 0.518 > \frac12 .
\]
Weddenschap twee: minstens één dubbele zes in $24$ worpen met twee
dobbelstenen,
\[
\P = 1 - \Bigl(\frac{35}{36}\Bigr)^{\!24} \approx 0.491 <
\frac12 .
\]
De Méré redeneerde dat $24$ worpen met kans $\frac1{36}$ zouden
moeten overeenkomen met $4$ worpen met kans $\frac16$ (dezelfde
verhouding $\frac{24}{36} = \frac46$); het falen van deze
evenredigheid --- kansen op verenigingen schalen niet lineair ---
zou zijn brief aan Pascal hebben uitgelokt, en daarmee de geboorte
van de kansrekening. De juiste vergelijking gaat via logaritmen:
$n$ pogingen met kans $p$ slagen minstens eenmaal met kans $1 -
(1-p)^n \approx 1 - \eu^{-np}$, dus is de eerlijke invariant $np$:
hier $4\cdot\frac16 = \frac23$ tegenover $24\cdot\frac1{36} =
\frac23$ --- gelijk! De twee weddenschappen verschillen pas op de
tweede orde in $p$, en net genoeg om er één over de vijftig-procent
lijn te duwen: kleine kansen vormen een gebied waar de intuïtie de
exponentiële functie nodig heeft, niet de liniaal.
\end{example}

\begin{remark}[Klassieke drogredenen bij het voorwaardelijk
stellen]
Drie terugkerende verwarringen, alle zichtbaar in de voorbeelden
hierboven. (i) \emph{Omkering}: $\pcond BA$ en $\pcond AB$
verschillen met de factor $\P(A)/\P(B)$ --- een test die $99\%$
nauwkeurig is op de zieken kan een positieve patiënt toch bijna
zeker gezond laten wanneer de ziekte zeldzaam is
(\cref{exo:b2:proba:3}); $\pcond{\text{ziek}}{ \text{positief}}$
noemen waar $\pcond{\text{positief}}{ \text{ziek}}$ wordt bedoeld,
is de drogreden van het basispercentage. (ii) \emph{Voorwaardelijk
stellen op de verkeerde \omterm{def:b2:proba:space}{gebeurtenis}}: bij Monty Hall is de juiste
voorwaardelijke \omterm{def:b2:proba:space}{gebeurtenis} ``de presentator opende deur $3$'',
niet ``de prijs zit niet achter deur $3$''; die twee dragen
verschillende informatie, en de hele $\frac23$ hangt van het
verschil af. (iii) \emph{Disjunct tegenover \omterm{def:b2:proba:independence}{onafhankelijk}}:
disjuncte \omterm{def:b2:proba:space}{gebeurtenissen} met positieve kans zijn nooit
\omterm{def:b2:proba:independence}{onafhankelijk} ($\P(A\cap B) = 0 \neq \P(A)\P(B)$) ---
\omterm{def:b2:proba:independence}{onafhankelijkheid} is verenigbaarheid van informatie, niet
afwezigheid van overlap.
\end{remark}

\begin{definition}[Onafhankelijkheid]\label{def:b2:proba:independence}
\omterm{def:b2:proba:space}{Gebeurtenissen} $A$ en $B$ heten
\emph{onafhankelijk}\index{onafhankelijke gebeurtenissen} wanneer
$\P(A \cap B) = \P(A)\P(B)$. Een familie $(A_i)_{i \in I}$
\omterm{def:b2:proba:space}{gebeurtenissen} heet \emph{(onderling) onafhankelijk} wanneer voor
elke eindige deelverzameling $J \subseteq I$ geldt
\[
\P\Bigl(\bigcap_{i \in J} A_i\Bigr)
= \prod_{i \in J} \P(A_i) .
\]
\end{definition}

\begin{remark}
Onderlinge \omterm{def:b2:proba:independence}{onafhankelijkheid} is strikt sterker dan paarsgewijze
\omterm{def:b2:proba:independence}{onafhankelijkheid}: bij twee eerlijke muntworpen zijn de
\omterm{def:b2:proba:space}{gebeurtenissen} ``de eerste is kop'', ``de tweede is kop'' en
``beide stemmen overeen'' paarsgewijs \omterm{def:b2:proba:independence}{onafhankelijk} (elk paar heeft
doorsnedekans $\frac14 = \frac12\cdot\frac12$), en toch heeft de
drievoudige doorsnede kans $\frac14 \neq \frac18$. Merk ook op dat
als $A, B$ \omterm{def:b2:proba:independence}{onafhankelijk} zijn, ook $A, B^c$ dat zijn (reken na:
$\P(A \cap B^c) = \P(A) - \P(A\cap B) = \P(A)(1 - \P(B))$), en dus
ook $A^c, B^c$.
\end{remark}

\begin{example}[Onafhankelijkheid afgelezen van een
productstructuur]\label{ex:b2:proba:productindep}
Gooi met twee eerlijke dobbelstenen: $\Omega = \intint16^2$ met
uniforme gewichten. Zij $A = {}$``eerste dobbelsteen even'' en $B =
{}$``tweede dobbelsteen minstens $5$''. Tellen: $\abs A = 3\cdot6 =
18$, $\abs B = 6\cdot2 = 12$, $\abs{A\cap B} = 3\cdot2 = 6$, dus
\[
\P(A\cap B) = \frac6{36} = \frac{18}{36}\cdot\frac{12}{36} =
\P(A)\,\P(B) :
\]
\omterm{def:b2:proba:independence}{onafhankelijk}, en het mechanisme is zichtbaar --- $A$ legt alleen
de eerste coördinaat vast, $B$ alleen de tweede, en de uniforme
maat op een productverzameling laat de coördinaattellingen
vermenigvuldigen. Elke bewering van het type ``\omterm{def:b2:proba:space}{gebeurtenissen} die
van disjuncte groepen worpen afhangen, zijn \omterm{def:b2:proba:independence}{onafhankelijk}''
(massaal gebruikt in de weekendopgave) is deze berekening, met meer
indices.
\end{example}

\begin{example}[Analyse van de eerste stap]\label{ex:b2:proba:firststep}
Wat is voor het meetkundige model van
\cref{ex:b2:proba:geometric} de kans $u$ dat de eerste kop op een
\emph{even} rang valt? Stel voorwaardelijk op de eerste worp: met
kans $p$ is de rang $1$ (oneven); met kans $q = 1 - p$ begint het
spel opnieuw met alle pariteiten omgekeerd, dus
\[
u = p\cdot0 + q\,(1 - u)
\qquad\Longrightarrow\qquad
u = \frac{q}{1 + q} .
\]
Eén regel, geen reeks --- en het stemt overeen met de
rechtstreekse sommatie van \cref{exo:b2:proba:9}, die $1 - u =
\frac1{1+q}$ geeft. Deze techniek van de ``eerste stap''
(voorwaardelijk stellen op het eerste experiment, en een verschoven
kopie van het probleem herkennen) is de kansrekenkundige vorm van
een recursie, en zij is de motor achter de vergelijkingen voor de
speelduur van \cref{exo:b2:proba:6} en de berekeningen van de
eerste doorgang in de weekendopgave.
\end{example}

\section{Het lemma van Borel--Cantelli}

\begin{definition}[Limes superior van gebeurtenissen]\label{def:b2:proba:limsup}
Voor een rij $(A_n)$ \omterm{def:b2:proba:space}{gebeurtenissen} is de \omterm{def:b2:proba:space}{gebeurtenis}
\[
\limsup_n A_n
= \bigcap_{N=0}^{\infty}\ \bigcup_{n \geq N} A_n
= \{\omega \in \Omega : \omega \in A_n
\text{ voor oneindig veel } n\}
\]
de \omterm{def:b2:proba:space}{gebeurtenis} ``$A_n$ treedt oneindig vaak op''.
\end{definition}

\begin{example}[``Oneindig vaak'' en ``uiteindelijk''
vertaald]\label{ex:b2:proba:limsuplang}
Het complement van $\limsup_nA_n$ is volgens de Morgan
\[
\Bigl(\bigcap_N\bigcup_{n\geq N}A_n\Bigr)^{\!c}
= \bigcup_N\bigcap_{n\geq N}A_n^c
= \{\omega : \omega \notin A_n \text{ voor alle grote }n\},
\]
de \omterm{def:b2:proba:space}{gebeurtenis} ``\emph{uiteindelijk} faalt $A_n$'' (geschreven
$\liminf_nA_n^c$). Dus zijn ``$A_n$ oneindig vaak'' en ``$A_n^c$
uiteindelijk'' complementair --- dit woordenboek recht houden
voorkomt de meeste ongelukken met kwantoren. Voorbeeldvertalingen
voor het muntwerpen: ``oneindig veel kop'' is $\limsup\{X_n =
H\}$; ``slechts eindig veel reeksen van $100$ maal kop'' is het
complement van een limsup; ``de lopende frequentie convergeert naar
$\frac12$'' is $\bigcap_j\bigcup_N\bigcap_{n\geq N}\{\abs{\widehat
p_n - \tfrac12} < \tfrac1j\}$ --- overal \omterm{def:b2:structures:countable}{aftelbare} bewerkingen,
zodat dit alle eerlijke \omterm{def:b2:proba:space}{gebeurtenissen} zijn.
\end{example}

\begin{theorem}[Borel--Cantelli]\label{thm:b2:proba:borelcantelli}
\leavevmode
\begin{enumerate}
\item Is $\sum_{n} \P(A_n) < \infty$, dan is
$\P\bigl(\limsup_n A_n\bigr) = 0$.
\item Zijn de \omterm{def:b2:proba:space}{gebeurtenissen} $A_n$ \omterm{def:b2:proba:independence}{onafhankelijk} en $\sum_n \P(A_n)
= \infty$, dan is $\P\bigl(\limsup_n A_n\bigr) = 1$.
\end{enumerate}
\end{theorem}

\begin{proof}
\emph{1.} Zij $C_N = \bigcup_{n \geq N}A_n$; de rij $(C_N)$ is
dalend met doorsnede $\limsup A_n$, en volgens de \omterm{def:b2:structures:countable}{aftelbare}
subadditiviteit (\cref{cor:b2:proba:subadd}) is
\[
\P(C_N) \leq \sum_{n \geq N}\P(A_n)
\xrightarrow[N\to\infty]{} 0
\]
(staart van een convergente reeks). De monotone \omterm{def:b2:metric:continuity}{continuïteit}
(\cref{thm:b2:proba:continuity}) besluit: $\P(\limsup A_n) = \lim_N
\P(C_N) = 0$.

\emph{2.} Het volstaat aan te tonen dat $\P\bigl(\bigcup_{n\geq
N}A_n\bigr) = 1$ voor elke $N$: hebben immers \omterm{def:b2:proba:space}{gebeurtenissen} $B_N$
alle kans $1$, dan is
\[
\P\Bigl(\Bigl(\bigcap_NB_N\Bigr)^{\!c}\Bigr)
= \P\Bigl(\bigcup_NB_N^c\Bigr)
\leq \sum_N\P(B_N^c) = 0
\]
wegens de \omterm{def:b2:structures:countable}{aftelbare} subadditiviteit
(\cref{cor:b2:proba:subadd}), zodat de \omterm{def:b2:structures:countable}{aftelbare} doorsnede
$\limsup A_n = \bigcap_N\bigcup_{n\geq N}A_n$ nog altijd kans $1$
heeft. Houd $N$ vast en beschouw voor $M > N$ het complement:
\[
\P\Bigl(\bigcap_{n=N}^{M} A_n^c\Bigr)
= \prod_{n=N}^{M}\bigl(1 - \P(A_n)\bigr)
\leq \prod_{n=N}^{M} e^{-\P(A_n)}
= \exp\Bigl(-\sum_{n=N}^M \P(A_n)\Bigr) ,
\]
met de \omterm{def:b2:proba:independence}{onafhankelijkheid} van de complementen en de
convexiteitsgrens $1 - x \leq e^{-x}$. Als $M \to \infty$ streeft
de exponent naar $-\infty$ wegens de divergentie van de reeks, dus
geeft de monotone \omterm{def:b2:metric:continuity}{continuïteit} (dalende rij) dat
$\P\bigl(\bigcap_{n \geq N}A_n^c\bigr) = 0$, dat wil zeggen
$\P\bigl(\bigcup_{n \geq N}A_n\bigr) = 1$.
\end{proof}

\begin{example}[Oneindige reeksen kop]\label{ex:b2:proba:runs}
Werp eeuwig met een eerlijke munt, en zij $A_n$ de \omterm{def:b2:proba:space}{gebeurtenis} ``de
worpen $n, n+1, \dots, n + k - 1$ zijn alle kop'' (een reeks van
$k$ maal kop die op tijdstip $n$ begint), voor vaste $k$. De
\omterm{def:b2:proba:space}{gebeurtenissen} $A_{jk}$ ($j = 1, 2, \dots$), die van disjuncte
blokken worpen afhangen, zijn \omterm{def:b2:proba:independence}{onafhankelijk}, elk met kans $2^{-k}$,
en $\sum_j 2^{-k} = \infty$: volgens Borel--Cantelli~2 zijn met
kans $1$ oneindig veel blokken \omterm{def:b2:metric:complete}{volledig} kop --- \emph{elk} vast
patroon keert bijna zeker oneindig vaak terug. Laten wij omgekeerd
de lengte van de reeks groeien, dan heeft $B_n = {}$``een reeks van
$2\log_2 n$ maal kop begint in $n$'' de kans $\P(B_n) = n^{-2}$,
\omterm{def:b2:series:summable}{sommeerbaar}, dus beginnen er bijna zeker slechts eindig veel van
zulke lange reeksen: Borel--Cantelli ijkt precies \emph{hoe lang}
de langste reeksen zijn.
\end{example}

\begin{example}[De oneindige aap, gekwantificeerd]
Een aap typt \omterm{def:b2:proba:independence}{onafhankelijke} uniforme letters uit een alfabet van
$26$ letters. Snijd het typoscript in disjuncte blokken van vier
letters; de \omterm{def:b2:proba:space}{gebeurtenissen} $A_j = {}$``blok $j$ spelt MATH'' zijn
\omterm{def:b2:proba:independence}{onafhankelijk} met $\P(A_j) = 26^{-4}$, en $\sum_j\P(A_j) = \infty$:
volgens Borel--Cantelli~2 typt de aap bijna zeker oneindig vaak
MATH --- en hetzelfde geldt voor elke vaste tekst van elke lengte,
met aangepaste blokken. De kwantitatieve voetnoot ontnuchtert het
wonder: $26^4 = 456\,976$, dus vergt de eerste MATH gemiddeld
ongeveer een half miljoen aanslagen, en een toneelstuk van
Shakespeare met $10^5$ tekens wacht op de orde van $26^{10^5}$
blokken --- bijna zeker is een uitspraak over de horizon $\infty$,
niet over enige horizon die een aap zal ontmoeten.
Borel--Cantelli bevestigt de limiet; de grootte van de termen
vertelt het verhaal op menselijke schaal.
\end{example}

\begin{remark}
In \cref{ex:b2:proba:runs} is de onderliggende \omterm{def:b2:proba:space}{uitkomstenruimte}
(oneindige rijen worpen) overaftelbaar, dus leeft het voorbeeld
strikt genomen in het maattheoretische kader van bachelorjaar 3; de
\emph{berekeningen} gebruiken echter alleen de regels die in dit
hoofdstuk zijn bewezen, toegepast op \omterm{def:b2:proba:space}{gebeurtenissen} die door
eindig veel worpen worden bepaald en op hun \omterm{def:b2:structures:countable}{aftelbare} combinaties.
Dat is de gebruikelijke afspraak op dit niveau: de theorie wordt op
\omterm{def:b2:structures:countable}{aftelbare} ruimten geformuleerd, en voorbeelden met oneindige spelen
worden met hetzelfde gereedschap behandeld.
\end{remark}

\begin{remark}[Vooruitblik binnen dit volume]
De machinerie van dit hoofdstuk wordt door de volgende twee in haar
geheel verbruikt. Indicatoren maken van \omterm{def:b2:proba:space}{gebeurtenissen}
toevalsveranderlijken, en de $\sigma$-additiviteit wordt de
\omterm{def:b2:series:summable}{sommeerbaarheid} die de verwachtingswaarde definieert
(\cref{ch:b2:randomvar}); Borel--Cantelli plus een \omterm{def:b2:series:summable}{sommeerbare}
staartgrens is precies hoe de sterke wet van de grote aantallen
voor munten daar wordt bewezen. In \cref{ch:b2:genfun} duikt de
monotone \omterm{def:b2:metric:continuity}{continuïteit} op het beslissende ogenblik weer op: de
uitstervingskans van een vertakkingsproces wordt \emph{gedefinieerd}
als de monotone limiet $\lim\P(Z_n = 0)$, en de
vastepuntsvergelijking waaraan zij voldoet, wordt verkregen door in
die stijgende rij naar de limiet over te gaan --- de laatste
stelling van het boek staat op de eerste stelling van dit
hoofdstuk.
\end{remark}

\begin{remark}[Methode: drie wegen naar kans één]
Uitspraken ``bijna zeker'' worden met drie hefbomen bewezen, in
stijgende volgorde van kracht. \emph{Monotone \omterm{def:b2:metric:continuity}{continuïteit}}: geef
de \omterm{def:b2:proba:space}{gebeurtenis} als een stijgende vereniging (of dalende doorsnede)
van \omterm{def:b2:proba:space}{gebeurtenissen} met eindige horizon en berekenbare kansen
(\cref{ex:b2:proba:sixeventually}). \emph{Nulverenigingen}: een
\omterm{def:b2:structures:countable}{aftelbare} vereniging van \omterm{def:b2:proba:space}{gebeurtenissen} met kans nul is nul
(\omterm{def:b2:structures:countable}{aftelbare} subadditiviteit), dus volstaat het elke slechte
\omterm{def:b2:proba:space}{gebeurtenis} afzonderlijk te doden --- zo voegt ``voor elke $j$ is
uiteindelijk $\abs{\widehat p_n - p} < 1/j$'' zich samen tot
convergentie. \emph{Borel--Cantelli}: is de \omterm{def:b2:proba:space}{gebeurtenis} een limsup,
sommeer dan de kansen; convergentie doodt haar (zonder
\omterm{def:b2:proba:independence}{onafhankelijkheid} nodig te hebben), en divergentie plus
\omterm{def:b2:proba:independence}{onafhankelijkheid} bevestigt haar. De juiste hefboom kiezen is
meestal het hele bewijs; de weekendopgave laat alle drie in één
argument lopen.
\end{remark}

\begin{remark}[Waar dit wordt gebruikt]
De monotone \omterm{def:b2:metric:continuity}{continuïteit} en Borel--Cantelli zijn de twee hefbomen
van elke uitspraak ``bijna zeker'': zij drijven de recurrentie van
de \omterm{pb:b2:proba:1}{toevalswandeling} in de weekendopgave van dit hoofdstuk aan, de
bijna-zekere kant van de wet van de grote aantallen
(\cref{ch:b2:randomvar}) en de analyse van het uitsterven van
vertakkingsprocessen (\cref{ch:b2:genfun}). Het volume van
bachelorjaar 3 bouwt de theorie opnieuw op met $\sigma$-algebra's
en de integraal van Lebesgue, waar de overaftelbare
\omterm{def:b2:proba:space}{uitkomstenruimten} die hier informeel worden gebruikt, \omterm{def:b2:metric:complete}{volledig}
streng worden.
\end{remark}

\section{Oefeningen}

\begin{exercise}[$\star$]\label{exo:b2:proba:1}
Een urne bevat $n$ genummerde ballen. De ballen worden één voor één
zonder teruglegging getrokken. Bereken de kans dat bal nummer $1$
vóór bal nummer $2$ wordt getrokken. Veralgemeen: de kans dat bal
$1$ als eerste wordt getrokken onder de ballen $1, \dots, k$.
\end{exercise}

\begin{exercise}[$\star$]\label{exo:b2:proba:2}
Toon aan dat de gewichten $p_k = \frac{1}{k(k+1)}$ op $\Omega =
\N^*$ een \omterm{def:b2:proba:space}{kansmaat} definiëren, en bereken $\P(2\N^*)$ (even
uitkomsten) als reeks; toon aan dat zij gelijk is aan $1 - \ln 2$.
\emph{(Telescopeer $\frac{1}{2j(2j+1)} = \frac{1}{2j} -
\frac{1}{2j+1}$ en gebruik de alternerende harmonische reeks,
\cref{ch:b2:series}.)}
\end{exercise}

\begin{exercise}[$\star$]\label{exo:b2:proba:3}
(Vals-positieven) Een ziekte treft één persoon op $10\,000$. Een
test spoort haar met kans $0.99$ op bij de zieken, en geeft met
kans $0.01$ een vals-positief bij de gezonden. Bereken de kans ziek
te zijn gegeven een positieve test, en geef commentaar.
\end{exercise}

\begin{exercise}[$\star\star$]\label{exo:b2:proba:4}
Zijn $A_1, \dots, A_n$ \omterm{def:b2:proba:space}{gebeurtenissen}. Bewijs de formule van de in-
en uitsluiting
\[
\P\Bigl(\bigcup_{i=1}^n A_i\Bigr)
= \sum_{\emptyset \neq J \subseteq \{1,\dots,n\}}
(-1)^{\abs J + 1}\,\P\Bigl(\bigcap_{i \in J}A_i\Bigr)
\]
door de identiteit $1 - \prod_{i=1}^n(1 - \mathbf{1}_{A_i}) =
\mathbf{1}_{\bigcup A_i}$ over $\Omega$ te integreren (dat wil
zeggen: te sommeren met gewichten $\P(\{\omega\})$).
\end{exercise}

\begin{exercise}[$\star\star$]\label{exo:b2:proba:5}
(Het probleem van de overeenkomsten, via in- en uitsluiting) $n$
brieven worden uniform willekeurig in $n$ enveloppen gestopt, één
per envelop. Toon met \cref{exo:b2:proba:4} aan dat de kans op
\emph{geen} enkele juiste overeenkomst gelijk is aan $\sum_{k=0}^n
\frac{(-1)^k}{k!} \to e^{-1}$, en leid de kans op precies één
overeenkomst af.
\end{exercise}

\begin{exercise}[$\star\star$]\label{exo:b2:proba:6}
Met een scheve munt (kanskop $p \in \intoo{0}{1}$) wordt geworpen
tot er tweemaal achter elkaar kop verschijnt. Zij $q_n$ de kans dat
het spel langer dan $n$ worpen duurt. Toon, door voorwaardelijk te
stellen op de eerste worp(en), aan dat $q_n = (1-p)\,q_{n-1} +
p(1-p)\,q_{n-2}$ voor $n \geq 2$, en leid af dat het spel met kans
$1$ eindigt. \emph{(Toon $q_n \to 0$ aan door met een meetkundige
rij te vergelijken: beide nulpunten van de karakteristieke
vergelijking hebben absolute waarde in $\intoo{0}{1}$.)}
\end{exercise}

\begin{exercise}[$\star\star\star$]\label{exo:b2:proba:7}
(Records) Trek een oneindige rij \omterm{def:b2:proba:independence}{onafhankelijke} uniforme
rangschikkingen, in de volgende combinatorische zin: voor elke $n$
is de onderlinge volgorde van de eerste $n$ trekkingen uniform over
de $n!$ mogelijkheden, en $R_n = {}$``de $n$-de trekking is een
record (groter dan alle vorige)''. Aangenomen dat de \omterm{def:b2:proba:space}{gebeurtenissen}
$R_n$ \omterm{def:b2:proba:independence}{onafhankelijk} zijn met $\P(R_n) = 1/n$ (bewijs minstens deze
laatste gelijkheid met de symmetrie), toon met Borel--Cantelli aan
dat er bijna zeker oneindig veel records optreden, maar dat records
op \emph{opeenvolgende} tijdstippen $n, n+1$ oneindig vaak optreden
met kans --- bereken $\sum_n \P(R_n \cap R_{n+1})$ en besluit wat
Borel--Cantelli~1 geeft.
\end{exercise}

\begin{exercise}[$\star\star\star$]\label{exo:b2:proba:8}
(In de stijl van Kochen--Stone, eenvoudiger versie) Zij $(A_n)$ een
rij \omterm{def:b2:proba:independence}{onafhankelijke gebeurtenissen} met $\P(A_n) = \frac{1}{n+1}$.
Toon aan dat $\P(\limsup A_n) = 1$, hoewel $\P(A_n) \to 0$:
``afzonderlijk zeldzaam, samen zeker''. Geef omgekeerd een rij
(afhankelijke) \omterm{def:b2:proba:space}{gebeurtenissen} met $\sum\P(A_n) = \infty$ en
$\P(\limsup A_n) = 0$, wat aantoont dat de \omterm{def:b2:proba:independence}{onafhankelijkheid} in
Borel--Cantelli~2 niet mag vervallen.
\end{exercise}

\begin{exercise}[$\star$]\label{exo:b2:proba:9}
Met een munt met kanskop $p \in \intoo01$ wordt geworpen tot de
eerste kop. Bereken de kans dat dit op een oneven rang gebeurt, en
evalueer haar voor een eerlijke munt.
\end{exercise}

\begin{exercise}[$\star\star$]\label{exo:b2:proba:10}
Zij $(A_n)_{n\geq1}$ een rij \omterm{def:b2:proba:independence}{onafhankelijke gebeurtenissen} met
$\P(A_n) = p_n < 1$. Toon aan dat
\[
\P\Bigl(\bigcap_{n\geq1}A_n^c\Bigr)
= \prod_{n\geq1}(1 - p_n)
:= \lim_{N\to\infty}\prod_{n=1}^N(1 - p_n),
\]
en dat deze limiet $> 0$ is dan en slechts dan als $\sum p_n <
\infty$. Breng dit in overeenstemming met Borel--Cantelli: is
$\sum p_n = \infty$, dan treedt niet alleen een zekere $A_n$ bijna
zeker op --- er treden er oneindig veel op.
\end{exercise}

\begin{exercise}[$\star\star$]\label{exo:b2:proba:11}
(De luciferdoosjes van Banach) Een roker houdt in elke zak een
doosje met $n$ lucifers en grijpt telkens in een uniform
willekeurige zak. Wanneer hij voor het eerst een doosje leeg
aantreft, wat is dan de kans dat het andere doosje precies $k$
lucifers bevat? Toon aan dat het antwoord
$\binom{2n-k}{n}2^{-(2n-k)}$ is en ga na dat deze kansen voor $n =
1$ tot $1$ sommeren.
\end{exercise}

\begin{exercise}[$\star\star\star$]\label{exo:b2:proba:12}
($\sigma$-additiviteit is een echt axioma)
(a) Toon aan dat er geen \omterm{def:b2:proba:space}{kansmaat} op $(\N, \mathcal P(\N))$ bestaat
die aan alle singletons hetzelfde gewicht geeft.
(b) Zij voor $A \subseteq \N^*$ de grootheid $d(A) =
\lim_n\frac{\abs{A\cap\intint1n}}{n}$ wanneer de limiet bestaat (de
\emph{natuurlijke dichtheid}). Toon aan dat $d$ eindig additief is
op paren waarvoor alle drie de dichtheden bestaan, dat zij aan elk
singleton dichtheid $0$ en aan $\N^*$ dichtheid $1$ geeft --- en
besluit dat $d$ niet $\sigma$-additief is.
(c) Geef een verzameling zonder dichtheid. \emph{(Neem de blokken
$\intint{2^{2k}}{2^{2k+1}-1}$ afwisselend wel en niet mee.)}
\end{exercise}

\section{Probleem: de eenvoudige toevalswandeling op $\Z$ is
recurrent}

\begin{omfigure}
\begin{tikzpicture}[scale=0.42, y=0.9cm]
  \draw[->] (0,0) -- (25,0) node[below] {$n$};
  \draw[->] (0,-2.6) -- (0,2.9) node[left] {$S_n$};
  \draw[omDef, thick] plot coordinates {(0,0) (1,1) (2,0)
    (3,-1) (4,0) (5,1) (6,2) (7,1) (8,2) (9,1) (10,0) (11,-1)
    (12,-2) (13,-1) (14,0) (15,-1) (16,0) (17,1) (18,2) (19,1)
    (20,0) (21,1) (22,0) (23,1) (24,2)};
  \foreach \t in {2,4,10,14,16,20,22}
    \fill[omThm] (\t,0) circle (5pt);
  \fill (0,0) circle (5pt);
\end{tikzpicture}

{\small Vierentwintig stappen van een \omterm{pb:b2:proba:1}{eenvoudige toevalswandeling};
de rode stippen markeren de terugkeren naar de oorsprong. De opgave
toont dat deze stippen met kans $1$ nooit ophouden te verschijnen
--- en toch heeft de wachttijd ertussen een divergent gemiddelde.}
\end{omfigure}

\begin{problem}[{Weekendopgave --- de recurrentiestelling van
Pólya op $\Z$, met onderweg het stemmenprobleem en een vleugje
arcsinus}]\label{pb:b2:proba:1}
Werp eeuwig met een eerlijke munt; zij $X_i = \pm1$ de $i$-de stap
en $S_n = X_1 + \dots + X_n$ de \emph{eenvoudige
toevalswandeling}\index{toevalswandeling} op $\Z$, met $S_0 = 0$.
Zoals in \cref{ex:b2:proba:runs} worden alle \omterm{def:b2:proba:space}{gebeurtenissen}
hieronder door eindig veel worpen bepaald of zijn zij \omterm{def:b2:structures:countable}{aftelbare}
combinaties van zulke \omterm{def:b2:proba:space}{gebeurtenissen}, en de \omterm{def:b2:proba:independence}{onafhankelijkheid} van
\omterm{def:b2:proba:space}{gebeurtenissen} die van disjuncte blokken worpen afhangen, hoort bij
het model. Wij schrijven $u_n = \P(S_{2n} = 0)$ en $N_n(k)$ voor
het aantal $\pm1$-paden van lengte $n$ van $0$ naar $k$.

\medskip
\noindent\textbf{Deel I --- Paden tellen.}
\begin{enumerate}
  \item Toon aan dat $N_n(k) = \binom{n}{(n+k)/2}$ wanneer $n + k$
        even is en $\abs k \leq n$, en $0$ anders; leid af dat
        $\P(S_n = k) = N_n(k)\,2^{-n}$. Waarom is elk afzonderlijk
        pad van lengte $n$ even waarschijnlijk?
  \item Toon aan dat $S_{2n+1} \neq 0$ en $u_n =
        \binom{2n}{n}4^{-n}$, en bereken $u_1, u_2, u_3$.
  \item Bewijs $u_n = \frac{2n-1}{2n}\,u_{n-1}$; leid af dat
        $(u_n)$ daalt naar $0$, en uit
        \cref{ex:b2:comparison:centralbinomial} dat
        \[
        u_n \sim \frac{1}{\sqrt{\pi n}},
        \qquad\text{zodat}\qquad
        \sum_n u_n = \infty .
        \]
  \item (Spiegelingsprincipe) Toon voor $k \geq 1$ aan dat de paden
        van lengte $n$ van $1$ naar $k$ die $0$ raken, in bijectie
        staan met de paden van $-1$ naar $k$; leid af dat het
        aantal paden van $0$ naar $k$ dat na tijdstip $0$ boven $0$
        blijft, gelijk is aan $N_{n-1}(k-1) - N_{n-1}(k+1)$.
  \item (Stemmenstelling) Leid af dat
        \[
        \P\bigl(S_1 > 0, \dots, S_{n-1} > 0 \bigm| S_n =
        k\bigr) = \frac kn \qquad (k \geq 1) :
        \]
        bij een telling waarin de winnaar met $k$ van de $n$
        stemmen voorstaat, is de kans dat hij gedurende de hele
        telling voorstond gelijk aan $k/n$. Ga dit met de hand na
        voor $n = 3$, $k = 1$.
\end{enumerate}

\medskip
\noindent\textbf{Deel II --- Terugkeer naar de oorsprong.}
\begin{enumerate}[resume]
  \item Bewijs de sleutelidentiteit
        \[
        \P(S_1 \neq 0,\ S_2 \neq 0,\ \dots,\ S_{2n} \neq 0) =
        u_n
        \]
        \emph{(stel voorwaardelijk op de eerste stap, sommeer de
        tellingen van vraag 4 over het eindpunt, en telescopeer;
        eindig met $2\binom{2n-1}{n} = \binom{2n}{n}$)}.
  \item Leid uit de monotone \omterm{def:b2:metric:continuity}{continuïteit}
        (\cref{thm:b2:proba:continuity}) af dat de wandeling met
        kans $1$ minstens eenmaal naar $0$ terugkeert, en dat $f_n
        := \P(\text{eerste terugkeer op tijdstip }2n)$ voldoet aan
        \[
        f_n = u_{n-1} - u_n = \frac{u_n}{2n-1},
        \qquad \sum_{n\geq1}f_n = 1 .
        \]
  \item Toon aan dat $\sum_n 2n\,f_n = \infty$: de terugkeer is
        zeker, maar de reeks die de gemiddelde wachttijd zou
        berekenen, divergeert (in de woordenschat van
        \cref{ch:b2:randomvar} heeft de terugkeertijd een oneindige
        verwachtingswaarde).
  \item Bewijs dat voor elke $k \geq 1$ geldt $\P(\text{minstens }
        k\text{ terugkeren naar }0) = 1$ \emph{(ontbind over de
        tijdstippen van de eerste $k$ terugkeren: de bijbehorende
        blokken worpen zijn disjunct, dus vermenigvuldigen de
        kansen en sommeren zij tot $(\sum_nf_n)^k$)}; besluit met
        de monotone \omterm{def:b2:metric:continuity}{continuïteit}:
        \[
        \P(S_n = 0 \text{ voor oneindig veel } n) = 1 :
        \]
        de \omterm{pb:b2:proba:1}{eenvoudige toevalswandeling} op $\Z$ is \emph{recurrent}.
  \item Toon aan dat de wandeling bijna zeker elke plaats $k \in
        \Z$ bezoekt, en dus (wegens de recurrentie, herstart bij
        het eerste bezoek) oneindig vaak. \emph{(De tekens van de
        opeenvolgende uitstapjes vanuit $0$ zijn \omterm{def:b2:proba:independence}{onafhankelijke}
        eerlijke munten; een positief uitstapje bezoekt $1$.)}
\end{enumerate}

\medskip
\noindent\textbf{Deel III --- Borel--Cantelli en de scheve
wandeling.}
\begin{enumerate}[resume]
  \item De \omterm{def:b2:proba:space}{gebeurtenissen} $A_n = \{S_{2n} = 0\}$ voldoen aan
        $\sum\P(A_n) = \infty$; leg uit waarom Borel--Cantelli~2
        er \emph{niet} op van toepassing is, en wat
        Borel--Cantelli~1 zou geven als de reeks convergeerde.
        (Dat is de strategie van het hele deel.)
  \item Laat de munt nu een scheefheid $p \neq \frac12$ hebben, $q
        = 1 - p$. Toon aan dat $\P(S_{2n} = 0) = \binom{2n}n(pq)^n
        = u_n\,(4pq)^n$ met $4pq < 1$, leid af dat
        $\sum_n\P(S_{2n} = 0) < \infty$, en besluit met
        Borel--Cantelli~1 dat de scheve wandeling bijna zeker
        slechts eindig vaak naar $0$ terugkeert.
  \item Nog steeds voor $p \neq \frac12$: toon aan dat $\P(S_n = k)
        \leq \binom{n}{\floor{n/2}}\,(pq)^{n/2}\,(p/q)^{k/2}$ voor
        elke vaste $k$, leid af dat elke plaats bijna zeker eindig
        vaak wordt bezocht, en besluit dat $\abs{S_n} \to \infty$
        bijna zeker: de scheve wandeling is \emph{transiënt}.
  \item Terug naar de eerlijke munt: bereken met vraag 6 de kans
        dat $200$ worpen \emph{geen} enkele gelijkstand opleveren
        ($S_n \neq 0$ voor $1 \leq n \leq 200$), numeriek
        $u_{100} \approx 0.056$. Geef commentaar op het trage
        verval $1/\sqrt{\pi n}$: gelijkstanden zijn op de lange
        duur zeker, maar zeldzamer dan de intuïtie doet vermoeden.
  \item (Eerste doorgang) Zij $T_1$ het eerste tijdstip waarop de
        wandeling $1$ raakt. Toon met het spiegelingsprincipe voor
        het maximum $M_n = \max_{i\leq n}S_i$ (bewezen in vraag 16,
        die niet van deze afhangt), of rechtstreeks uit vraag 7
        door voorwaardelijk op de eerste stap te stellen, aan dat
        $\P(T_1 = 2n - 1) = f_n$; leid af dat $\P(T_1 < \infty) =
        1$, terwijl de reeks $\sum(2n-1)f_n$ voor de gemiddelde
        tijd divergeert.
\end{enumerate}

\medskip
\noindent\textbf{Deel IV --- Maxima, laatste nulpunt, lange
voorsprongen.}
\begin{enumerate}[resume]
  \item (Spiegeling voor het maximum) Bewijs voor $k \geq 1$ dat
        \[
        \P(M_n \geq k) = 2\,\P(S_n > k) + \P(S_n = k)
        \]
        door het pad na zijn eerste bezoek aan niveau $k$ te
        spiegelen.
  \item Leid af dat $\P(M_{2n} \geq 1) = 1 - u_n$, dat wil zeggen
        $\P(S_i \leq 0 \text{ voor alle } i \leq 2n) = u_n$: de
        kans om nooit voor te staan is gelijk aan de kans om nooit
        op nul te zijn (vraag 6) --- twee verschillende
        \omterm{def:b2:proba:space}{gebeurtenissen}, één kans.
  \item (Laatste nulpunt) Zij $L_{2n} = \max\{k \leq 2n : S_k =
        0\}$ (even). Toon, door vraag 6 met de \omterm{def:b2:proba:independence}{onafhankelijkheid}
        van disjuncte blokken worpen te combineren, aan dat
        \[
        \P(L_{2n} = 2k) = u_k\,u_{n-k}
        \qquad (0 \leq k \leq n),
        \]
        en leid zonder verdere berekening de binomiale identiteit
        $\sum_{k=0}^n u_ku_{n-k} = 1$ af.
  \item Toon aan dat de verdeling van $L_{2n}$ \omterm{def:b2:quadratic:adjoint}{symmetrisch} is
        ($\P(L = 2k) = \P(L = 2n - 2k)$) en, met $u_j \sim
        1/\sqrt{\pi j}$, dat haar uitersten haar meest
        waarschijnlijke waarden zijn. Zet dit voor $n = 5$ in een
        tabel: $\P(L_{10} = 0) = u_5 \approx 0.246$ tegenover
        $\P(L_{10} = 4) = u_2u_3 \approx 0.117$. Interpreteer: in
        een lang eerlijk spel valt de laatste gelijkstand meestal
        heel vroeg of heel laat --- lange voorsprongen zijn de
        regel, niet de uitzondering.
  \item Zet de vragen 16--19 in elkaar tot een alinea over het
        beeld van de fluctuaties van de eerlijke wandeling: de
        diffusieve schaal die vraag 3 suggereert, de zekerheid van
        de terugkeer tegenover de divergente gemiddelde wachttijd,
        en de volharding van voorsprongen met haar arcsinusaroma.
\end{enumerate}

\medskip
\noindent\textbf{Deel V --- De vernieuwingsidentiteit en de
stelling van Pólya.}
\begin{enumerate}[resume]
  \item Bewijs, door $\{S_{2n} = 0\}$ te partitioneren over het
        tijdstip van de eerste terugkeer, de
        \emph{vernieuwingsidentiteit}
        \[
        u_n = \sum_{k=1}^n f_k\,u_{n-k} \quad (n \geq 1),
        \qquad\text{en dus}\qquad
        U(x)\bigl(1 - F(x)\bigr) = 1 \quad (0 \leq x < 1),
        \]
        waarbij $U(x) = \sum_{n\geq0}u_nx^n$ en $F(x) =
        \sum_{n\geq1}f_nx^n$ (verantwoord de stralen en het product
        van de reeksen met \cref{ch:b2:powerseries}).
  \item Leid de \emph{recurrentietweedeling} af: met $x \to 1^-$
        (monotone limieten van reeksen met niet-negatieve
        coëfficiënten) is
        \[
        \sum_n u_n = \infty \iff \sum_n f_n = 1 ,
        \]
        en toets haar aan de vragen 3, 7 (eerlijke wandeling) en 12
        (scheve wandeling).
  \item (Dimensie $2$) De eenvoudige wandeling op $\Z^2$ neemt
        uniform stappen $(\pm1, 0)$, $(0, \pm1)$. Toon aan dat de
        gedraaide coördinaten $U_n = X_n + Y_n$ en $V_n = X_n -
        Y_n$ \emph{\omterm{def:b2:proba:independence}{onafhankelijke}} eerlijke wandelingen op $\Z$
        uitvoeren, leid af dat
        \[
        \P\bigl(S^{(2)}_{2n} = (0,0)\bigr) = u_n^2 \sim
        \frac1{\pi n},
        \qquad \sum_n u_n^2 = \infty ,
        \]
        en besluit met de vragen 21--22 (waarvan de bewijzen
        woordelijk overgaan) dat de wandeling op $\Z^2$ recurrent
        is.
  \item (Dimensie $3$) Neem voor de eenvoudige wandeling op $\Z^3$
        de lokale schatting $\P(S^{(3)}_{2n} = 0) \leq
        C\,n^{-3/2}$ aan (bewezen met de lokale limietstelling in
        het volume van bachelorjaar 3). Leid met
        Borel--Cantelli~1 af dat de wandeling op $\Z^3$ transiënt
        is, en formuleer het volledige resultaat: de \emph{stelling
        van Pólya} --- de \omterm{pb:b2:proba:1}{eenvoudige toevalswandeling} is recurrent
        in de dimensies $1$ en $2$, en transiënt in dimensie $3$ en
        hoger.
  \item Synthese. Som de exacte rol op van: het tellen van paden en
        de spiegeling; de monotone \omterm{def:b2:metric:continuity}{continuïteit}; de
        \omterm{def:b2:proba:independence}{onafhankelijkheid} van disjuncte blokken worpen;
        Borel--Cantelli~1; de vernieuwingsidentiteit. Welk enkel
        \omterm{def:b2:powerseries:analytic}{analytisch} feit ($u_n \sim 1/\sqrt{\pi n}$, dus $\sum u_n =
        \infty$ maar ook $\sum u_n^2 = \infty$ en $\sum n^{-3/2} <
        \infty$) beslist in elke dimensie tussen recurrentie en
        transiëntie?
\end{enumerate}
\end{problem}
